%%%%%%%%%%%%%%%%%%%%%%%%%%%%%%%%%%%%%%%%%%%%%%%%%%%%%%%%%%%%%%%%%%%%%%%
%%%%%%%%%%%%%%%%%%%%%% Simple LaTeX CV Template %%%%%%%%%%%%%%%%%%%%%%%%
%%%%%%%%%%%%%%%%%%%%%%%%%%%%%%%%%%%%%%%%%%%%%%%%%%%%%%%%%%%%%%%%%%%%%%%%

%%%%%%%%%%%%%%%%%%%%%%%%%%%%%%%%%%%%%%%%%%%%%%%%%%%%%%%%%%%%%%%%%%%%%%%%
%% NOTE: If you find that it says                                     %%
%%                                                                    %%
%%                           1 of ??                                  %%
%%                                                                    %%
%% at the bottom of your first page, this means that the AUX file     %%
%% was not available when you ran LaTeX on this source. Simply RERUN  %%
%% LaTeX to get the ``??'' replaced with the number of the last page  %%
%% of the document. The AUX file will be generated on the first run   %%
%% of LaTeX and used on the second run to fill in all of the          %%
%% references.                                                        %%
%%%%%%%%%%%%%%%%%%%%%%%%%%%%%%%%%%%%%%%%%%%%%%%%%%%%%%%%%%%%%%%%%%%%%%%%

%%%%%%%%%%%%%%%%%%%%%%%%%%%% Document Setup %%%%%%%%%%%%%%%%%%%%%%%%%%%%

% Don't like 10pt? Try 11pt or 12pt
\documentclass[10pt]{article}

% This is a helpful package that puts math inside length specifications
\usepackage{calc}


% Simpler bibsection for CV sections
% (thanks to natbib for inspiration)
\makeatletter
\newlength{\bibhang}
\setlength{\bibhang}{1em}
\newlength{\bibsep}
 {\@listi \global\bibsep\itemsep \global\advance\bibsep by\parsep}
\newenvironment{bibsection}%
        {\vspace{-\baselineskip}\begin{list}{}{%
       \setlength{\leftmargin}{\bibhang}%
       \setlength{\itemindent}{-\leftmargin}%
       \setlength{\itemsep}{\bibsep}%
       \setlength{\parsep}{\z@}%
        \setlength{\partopsep}{0pt}%
        \setlength{\topsep}{0pt}}}
        {\end{list}\vspace{-.6\baselineskip}}
\makeatother

% Layout: Puts the section titles on left side of page
\reversemarginpar

%
%         PAPER SIZE, PAGE NUMBER, AND DOCUMENT LAYOUT NOTES:
%
% The next \usepackage line changes the layout for CV style section
% headings as marginal notes. It also sets up the paper size as either
% letter or A4. By default, letter was used. If A4 paper is desired,
% comment out the letterpaper lines and uncomment the a4paper lines.
%
% As you can see, the margin widths and section title widths can be
% easily adjusted.
%
% ALSO: Notice that the includefoot option can be commented OUT in order
% to put the PAGE NUMBER *IN* the bottom margin. This will make the
% effective text area larger.
%
% IF YOU WISH TO REMOVE THE ``of LASTPAGE'' next to each page number,
% see the note about the +LP and -LP lines below. Comment out the +LP
% and uncomment the -LP.
%
% IF YOU WISH TO REMOVE PAGE NUMBERS, be sure that the includefoot line
% is uncommented and ALSO uncomment the \pagestyle{empty} a few lines
% below.
%

%% Use these lines for letter-sized paper
\usepackage[paper=letterpaper,
            includefoot, % Uncomment to put page number above margin
            marginparwidth=1in,     % Length of section titles
            marginparsep=.08in,       % Space between titles and text
            margin=0.5in,            % 0.75 inch margins
            includemp]{geometry}

%% Use these lines for A4-sized paper
%\usepackage[paper=a4paper,
%            %includefoot, % Uncomment to put page number above margin
%            marginparwidth=30.5mm,    % Length of section titles
%            marginparsep=1.5mm,       % Space between titles and text
%            margin=25mm,              % 25mm margins
%            includemp]{geometry}

%% More layout: Get rid of indenting throughout entire document
\setlength{\parindent}{0in}

%% This gives us fun enumeration environments. compactitem will be nice.
\usepackage{paralist}

%% Reference the last page in the page number
%
% NOTE: comment the +LP line and uncomment the -LP line to have page
%       numbers without the ``of ##'' last page reference)
%
% NOTE: uncomment the \pagestyle{empty} line to get rid of all page
%       numbers (make sure includefoot is commented out above)
%
\usepackage{fancyhdr,lastpage}
\pagestyle{fancy}
\pagestyle{empty}      % Uncomment this to get rid of page numbers
\fancyhf{}\renewcommand{\headrulewidth}{0pt}
\fancyfootoffset{\marginparsep+\marginparwidth}
\newlength{\footpageshift}
\setlength{\footpageshift}
          {0.5\textwidth+0.5\marginparsep+0.5\marginparwidth-2in}
\lfoot{\hspace{\footpageshift}%
       \parbox{4in}{\, \hfill %
                    \arabic{page} of \protect\pageref*{LastPage} % +LP
%                    \arabic{page}                               % -LP
                    \hfill \,}}

% Finally, give us PDF bookmarks
\usepackage{color,hyperref}
\definecolor{darkblue}{rgb}{0.0,0.0,0.55}
\hypersetup{colorlinks,breaklinks,
            linkcolor=darkblue,urlcolor=darkblue,
            anchorcolor=darkblue,citecolor=darkblue}

%%%%%%%%%%%%%%%%%%%%%%%% End Document Setup %%%%%%%%%%%%%%%%%%%%%%%%%%%%


%%%%%%%%%%%%%%%%%%%%%%%%%%% Helper Commands %%%%%%%%%%%%%%%%%%%%%%%%%%%%

% The title (name) with a horizontal rule under it
%
% Usage: \makeheading{name}
%
% Place at top of document. It should be the first thing.
\newcommand{\makeheading}[1]%
        {\hspace*{-\marginparsep minus \marginparwidth}%
         \begin{minipage}[t]{\textwidth+\marginparwidth+\marginparsep}%
                {\large \bfseries #1}\\[-0.15\baselineskip]%
                 \rule{\columnwidth}{1.5pt}%
         \end{minipage}}

% The section headings
%
% Usage: \section{section name}
%
% Follow this section IMMEDIATELY with the first line of the section
% text. Do not put whitespace in between. That is, do this:
%
%       \section{My Information}
%       Here is my information.
%
% and NOT this:
%
%       \section{My Information}
%
%       Here is my information.
%
% Otherwise the top of the section header will not line up with the top
% of the section. Of course, using a single comment character (%) on
% empty lines allows for the function of the first example with the
% readability of the second example.
\renewcommand{\section}[2]%
        {\pagebreak[2]\vspace{1\baselineskip}%
         \phantomsection\addcontentsline{toc}{section}{#1}%
         \hspace{0in}%
         \marginpar{
         \raggedright \scshape #1}#2}

% An itemize-style list with lots of space between items
\newenvironment{outerlist}[1][\enskip\textbullet]%
        {\begin{itemize}[#1]}{\end{itemize}%
         \vspace{-0.6\baselineskip}}

% An environment IDENTICAL to outerlist that has better pre-list spacing
% when used as the first thing in a \section
\newenvironment{lonelist}[1][\enskip\textbullet]%
        {\vspace{-\baselineskip}\begin{list}{#1}{%
        \setlength{\partopsep}{0pt}%
        \setlength{\topsep}{0pt}}}
        {\end{list}\vspace{-.6\baselineskip}}

% An itemize-style list with little space between items
\newenvironment{innerlist}[1][\enskip\textbullet]%
        {\begin{compactitem}[#1]}{\end{compactitem}}

% An environment IDENTICAL to innerlist that has better pre-list spacing
% when used as the first thing in a \section
\newenvironment{loneinnerlist}[1][\enskip\textbullet]%
        {\vspace{-\baselineskip}\begin{compactitem}[#1]}
        {\end{compactitem}\vspace{-.6\baselineskip}}

% To add some paragraph space between lines.
% This also tells LaTeX to preferably break a page on one of these gaps
% if there is a needed pagebreak nearby.
\newcommand{\blankline}{\vspace*{0pt}}

% Uses hyperref to link DOI
\newcommand\doilink[1]{\href{http://dx.doi.org/#1}{#1}}
\newcommand\doi[1]{doi:\doilink{#1}}


%%%%%%%%%%%%%%%%%%%%%%%% End Helper Commands %%%%%%%%%%%%%%%%%%%%%%%%%%%

%%%%%%%%%%%%%%%%%%%%%%%%% Begin CV Document %%%%%%%%%%%%%%%%%%%%%%%%%%%%

\usepackage{graphicx}
\begin{document}

%\href{http://google.com}{\includegraphics[bb=-4 30 0 0,scale=0.7]{git.png}}


% For a right forkme png
%\begin{figure}%[4]
%		\vspace{-2.6cm}
%		\hspace{12.2cm}
%		\href{https://github.com/tUrG0n/Resume}{\includegraphics[scale=0.7]{git.png}}
%		\vspace{-2.0cm}
%\end{figure}


% For a left forkme png
%\begin{figure}%[4]
%    \flushleft
%		\vspace{-2.6cm}
%		\hspace{-6.0cm}
%		\href{https://github.com/tUrG0n/Resume}{\includegraphics[scale=0.7]{git.png}}
%		\vspace{-2.0cm}
%\end{figure}


\makeheading{Karan Jitendra Thakkar \hfill \href{http://in.linkedin.com/in/karanjthakkar}{Linkedin}}

\section{Contact}
%
% NOTE: Mind where the & separators and \\ breaks are in the following
%       table.
%
% ALSO: \rcollength is the width of the right column of the table
%       (adjust it to your liking; default is 1.85in).
%
\newlength{\rcollength}\setlength{\rcollength}{1.65in}%
\vspace{0.1cm}
\begin{tabular}[t]{@{}c|c|c}
\textit{Email:} \href{mailto:karanjthakkar@gmail.com}{karanjthakkar@gmail.com} & \textit{Website:} \href{karanjthakkar.co.nr}{karanjthakkar.co.nr} & \textit{Mobile:} $+$91-9637799260 \\

\end{tabular}

\hspace{-2.85cm}
\rule{1.185\linewidth}{1pt}%
\blankline

\section{Summary}

\vspace{-10pt}

\begin{innerlist}
\item Hybrid Mobile Application Developer with 1.5 years of experience
\item Comfortable with Grunt for seamless build tooling
\item Capable of handling challenging problems and come up with creative solutions
\item Love writing code that follows best practices, is maintainable and awesome!
\end{innerlist}

\blankline

\section{Education}
\href{http://kitcoek.in/}{\textbf{KIT's College of Engineering}},
Kolhapur, India
\begin{innerlist}
\item[] B.E., \href{http://kitcoek.in/academics/electronics-telecommunications/} {Electronics and Telecommunication Engineering}
\hfill \textbf{August '08 - July '12} 
\item[] \textit{First Class with Distinction}, Average grade 70.68\%
\end{innerlist}

\blankline

\section{Online Courses}
\href{http://education.mongodb.com/}{\textbf{M101JS: MongoDB for Node.js Developers}} \small \hfill \textbf{November '13 - December '13}
\begin{innerlist}
\item[] Grade 82.00\% \hfill \small \href{http://education.mongodb.com/downloads/certificates/2bc6c92e90974667bc5e5c93db335e71/Certificate.pdf}{\emph{Completion Certificate}}
\end{innerlist}

\href{http://coursera.org/}{\textbf{Computing for Data Analysis}} \small \hfill \textbf{October '13 - November '13}
\begin{innerlist}
\item[] Grade 89.00\% \hfill \small \href{http://karanjthakkar.files.wordpress.com/2013/11/computing-for-data-analysis.pdf}{\emph{Statement of Accomplishment}}
\end{innerlist}

\href{http://ai-class.com/}{\textbf{Introduction to Artificial Intelligence}} \small \hfill \textbf{October '12 - December '12}
\begin{innerlist}
\item[] Grade 79.30\% \hfill \small \href{http://karanjthakkar.files.wordpress.com/2012/12/statement-of-accomplishment.pdf}{\emph{Statement of Accomplishment}}
\end{innerlist}

\blankline

\section{Tech Skills}
\textbf{Languages:} Javascript, HTML5, CSS3, C$+$$+$, C$\#$, Core Java

\vspace{8pt}

\textbf{Tools:} Brackets, Eclipse, TortoiseSVN, Visual Studio, git, JIRA, HP Quality Center, \LaTeX

\vspace{8pt}

\textbf{Frameworks:} BackboneJS, Sencha Touch, EnyoJS

\vspace{8pt}

\textbf{Libraries:} Phonegap, jQuery, Underscore, RequireJS, Grunt

\blankline

\section{Professional Experience}
\small \textbf{\href{http://www.tcs.com}{Tata Consultancy Services Ltd.}}, Pune, India \normalsize
\begin{innerlist}
\item[] \textbf{Assistant Systems Engineer \small \hfill January '13 - Present}
\item[] Worked on the development of a hybrid mobile applications using HTML, CSS, Javascript, Phonegap and modular cross platform javascript frameworks like BackboneJS, Sencha Touch, EnyoJS
\end{innerlist}

\vspace{8pt}

\small \textbf{\href{http://www.autonomous-systems.in/}{Gade Autonomous Systems Pvt. Ltd.}}, Mumbai, India \normalsize
\begin{innerlist}
\item[] \textbf{Research Intern \small \hfill September '12 - November '12}
\item[] Developed a WPF application for a Gesture Based Picture Viewer using Microsoft Kinect
\end{innerlist}

\vspace{8pt}

\small \textbf{\href{https://www.linkedin.com/company/callbackers-pvt-ltd}{Callbackers Pvt. Ltd.}}, Chandigasrh, India \normalsize
\begin{innerlist}
\item[] \textbf{Remote Intern \small \hfill July '12 - August '12}
\item[] Designed ready-to-manufacture GSM and Xbee shields for Arduino Uno using EAGLE PCB design tool
\end{innerlist}

\blankline

\section{Projects}
\small \textbf{Investment Pulse \hfill December '13 - Present} \normalsize
\begin{innerlist}
\item[] \textbf{Client}: Fidelity Worldwide Investments
\item[] \textbf{Technologies}: HTML5, CSS3, Javascript, BackboneJS, Phonegap, Google Analytics, iOS/Android Plugins
\item[] Developed an application that allowed individual investors to view details of their investments, related news and manage their investments \normalsize
\end{innerlist}

\vspace{8pt}

\small \textbf{Resource Management System \hfill August '13 - October '13} \normalsize
\begin{innerlist}
\item[] \textbf{Client}: TCS Internal - Proof of Concept
\item[] \textbf{Technologies}: HTML5, CSS3, Javascript, Sencha Touch, ExtJS
\item Front End Development utilizing HTML5, CSS3, and JavaScript leveraging the Sencha Touch cross platform framework and a Java/MySQL backend
\hspace{-2.85cm}\item UI/UX Development with responsive web design principles \normalsize
\end{innerlist}

\vspace{8pt}

\small \textbf{HP\textsuperscript{\textregistered} Discover \hfill December '13 - Present} \normalsize
\begin{innerlist}
\item[] \textbf{Client}: Hewlett-Packard
\item[] \textbf{Technologies}: HTML5, CSS3, Javascript, EnyoJS
\item[] Worked in a limited capacity for enhancement of an event mangement application for attendees of HP\textsuperscript{\textregistered} Discover 2013 \normalsize
\end{innerlist}

\vspace{8pt}

\small \textbf{V'Me: A Gesture Based Picture Viewer application \hfill September '12 - November '12} \normalsize
\begin{innerlist}
\item[] \textbf{Client}: Gade Autonomous Systems Pvt. Ltd.
\item[] \textbf{Technologies}: C\#, XAML, Kinect SDK
\item[] Developed a WPF application for a Gesture Based Picture Viewer using Microsoft Kinect. Implemented functionalities such as: Swipe Left, Swipe Right, Zoom In, Zoom Out, Panning, Moving mouse cursor and Switching to different folders.\hfill \small \href{http://github.com/karanjthakkar/Gesture-Based-Picture-Viewer}{\textit{Code}} $|$ \href{http://youtu.be/v8SumS-I1qo}{\textit{Demo}} $|$ \href{http://autonomous-systems.in/#vme}{\textit{Website}} \normalsize
\end{innerlist}

\vspace{8pt}

Other independent projects hosted on \href{http://github.com/karanjthakkar}{\textit{\textbf{Github} } }

\blankline


\end{document}
%%%%%%%%%%%%%%%%%%%%%%%%%% End CV Document %%%%%%%%%%%%%%%%%%%%%%%%%%%%%

